\section{Automated Red Teaming}\label{sec:automated_red_teaming}

\subsection{Problem Definition and Metrics}\label{sec:problem_definition}
Following \citet{perez2022red, zou2023universal}, we formulate the red teaming task as designing test cases $\{x_1, x_2, \dotsc, x_N\}$ in order to elicit a given behavior $y$ from one or more target LLMs.

The primary measure of a red teaming method's success is its attack success rate (ASR) on a given target model, which is the percentage of test cases that elicit the behavior from the target model. To improve evaluation efficiency, we follow prior work in assuming that target models generate completions deterministically using greedy decoding \citep{zou2023universal, chao2023jailbreaking, mehrotra2023treeOfAttacks}. Formally, let $f$ be a target model with generation function $f_T(x) = x'$, where $T$ is the number of tokens to be generated, $x$ is a test case, and $x'$ is the completion. Let $g$ be a red teaming method that generates a list of test cases, and let $c$ be a classifier mapping completion $x'$ and behavior $y$ to $1$ if a test case was successful and $0$ if not. The ASR of $g$ on target model $f$ for behavior $y$ is then defined as
\[\text{ASR}(y, g, f) = \frac{1}{N}\sum c(f_T(x_i), y).\]



\subsection{Toward Improved Evaluations}\label{sec:improved_evaluations}
In prior work, a range of specialized evaluation setups have been proposed. However, there has not yet been a systematic effort to standardize these evaluations and provide a large-scale comparison of existing methods. Here, we discuss key qualities for automated red teaming evaluations, how existing evaluations fall short, and how we improve on them. Namely, we identify three key qualities: breadth, comparability, and robust metrics.

\paragraph{Breadth.}
Red teaming methods that can only obtain high ASR on a small set of harmful behaviors may be less useful in practice. Thus, it is desirable for evaluations to encompass a wide variety of harmful behaviors. We conduct a brief survey of prior evaluations to tabulate their diversity of behaviors, finding that most use short, unimodal behaviors with less than $100$ unique behaviors. By contrast, our proposed benchmark contains novel functional categories and modalities of behaviors, including contextual behaviors, copyright behaviors, and multimodal behaviors. These results are shown in \Cref{table:dataset_comparision}.

In addition to providing comprehensive tests of red teaming methods, a broad array of behaviors can greatly enhance the evaluation of defenses. Different developers may be interested in preventing different behaviors. Our benchmark includes a broad range of behavior categories to enable tailored evaluations of defenses. Moreover, we propose a standardized evaluation framework that can be easily extended to include new undesired behaviors, allowing developers to quickly evaluate their defenses against a wide range of attacks on the behaviors that are most concerning to them.


\paragraph{Comparability.}
The foundation of any evaluation is being able to meaningfully compare different methods. One might naively think that running off-the-shelf code is sufficient for comparing the performance of one's red teaming method to a baseline. However, we find that there is considerable nuance to obtaining a fair comparison. In particular, we identify a crucial factor overlooked in prior work that highlights the importance of standardization.


In developing our evaluations, we found that the number of tokens generated during evaluation can have a drastic effect on ASR computed by substring matching metrics used in earlier work. We show in \Cref{fig:num_tokens} that the choice of this parameter can change ASR by up to $30\%$. Unfortunately, this parameter has not been standardized in prior work, rendering cross-paper comparisons effectively meaningless. In \Cref{sec:evaluation_pipeline}, we propose a new metric that is more robust to variations in this parameter, and we standardize this parameter to $N=512$ to allow the metric to converge.


\begin{figure}[t]
    \centering
    \includegraphics[width=0.49\textwidth]{figures/num_tokens.pdf}
    \vspace{-20pt}
    \caption{The number of tokens generated by the target model during evaluation drastically impacts the attack success rate (ASR) of red teaming methods. This crucial evaluation parameter is not standardized in prior work. As a result, cross-paper comparisons can be misleading.}
    \label{fig:num_tokens}
    \vspace{-10pt}
\end{figure}

\paragraph{Robust Metrics.}
Research on red teaming LLMs benefits from the codevelopment of attacks and defenses. However, this means that metrics for evaluating red teaming methods can face considerable optimization pressure as both attacks and defenses seek to improve performance. As a result, one cannot simply use any classifier for this process. As a prequalification, classifiers should exhibit robustness to nonstandard scenarios, lest they be easily gamed. Here, we propose an initial prequalification test consisting of three types of nonstandard test case completions:
\begin{enumerate}
    \item Completions where the model initially refuses, but then continues to exhibit the behavior
    \item Random benign paragraphs
    \item Completions for unrelated harmful behaviors
\end{enumerate}
We compare a variety of classifiers on these sets in \Cref{table:cls_results_sets}, finding that many previously used classifiers lack robustness to these simple but nonstandard scenarios. Additionally, a crucial measure to ensuring the robustness of evaluation metrics is using held-out classifiers and a validation/test split for harmful behaviors. We find that several prior works directly evaluate on the metric optimized by their method---a practice that can lead to substantial gaming.


